% Generated draft: requires mathematical and source review.
\documentclass[11pt]{article}
\usepackage[margin=25mm]{geometry}
\usepackage{fontspec}
\setmainfont{Latin Modern Roman}
\usepackage{amsmath,amssymb,mathtools,unicode-math}
\setmathfont{Latin Modern Math}
\usepackage{enumitem,needspace,xurl,hyperref}
\setlength{\parindent}{0pt}
\setlength{\parskip}{5pt}
\setlength{\emergencystretch}{3em}
\newcommand{\N}{\mathbb{N}}
\newcommand{\Z}{\mathbb{Z}}
\newcommand{\Q}{\mathbb{Q}}
\newcommand{\R}{\mathbb{R}}
\newcommand{\C}{\mathbb{C}}
\newcommand{\F}{\mathbb{F}}
\newcommand{\A}{\mathbb{A}}
\newcommand{\PP}{\mathbb{P}}
\newcommand{\E}{\mathbb{E}}
\newcommand{\eps}{\varepsilon}
\newcommand{\dd}{\,\mathrm{d}}
\begin{document}
\section*{UnsolvedMath ID: 30005267}
\textbf{Source problem code:} catalogue entry
\subsection{Definitions and mathematical statement}
Let \(\mathbb P=\{p\in\Z:p>1\text{ and the only positive divisors of }p\text{ are }1,p\}\). Let \(P\in\Z[X]\) be irreducible in \(\Z[X]\), with \(\deg P\geq2\). Here irreducibility means that \(P\) cannot be expressed as a product of two nonunits of \(\Z[X]\); its units are \(1\) and \(-1\). Say that \(P\) has no fixed prime divisor if
\[
\forall p\in\mathbb P\;\exists m\in\Z:\quad p\nmid P(m).
\]
The supplied assertion is that every such polynomial satisfies
\[
\#\{n\in\Z:P(n)\in\mathbb P\}=\infty.
\]
\textbf{Unresolved sign qualification.} The supplied statement imposes no sign condition on \(P\) and does not replace \(P(n)\) by \(\lvert P(n)\rvert\). With positive primes as defined above, polynomials of even degree with negative leading coefficient are eventually negative in both integer directions, so the displayed assertion cannot hold for such polynomials. The intended formulation therefore requires clarification: a positivity hypothesis or an absolute-value convention would change the supplied assertion. Neither is silently imposed here. The domain \(n\in\Z\) is retained exactly as supplied.

The original extraction also contains adjacent material on largest prime factors of quartic polynomial values in the cyclic and dihedral cases, attributed to C\'ecile Dartyge, in joint work with James Maynard. It attributes a general bound to Tenenbaum, reference [6], under the conditions that \(f\in\Z[X]\) is irreducible, \(\deg f\geq2\), \(0<\alpha<2-\log4\), and \(x\geq x_0(\alpha,f)\). The extracted product and largest-prime-factor notation are corrupted, so no precise auxiliary inequality is asserted here.
\subsection{Short English statement}
The supplied question asks whether every irreducible integer polynomial of degree at least two, with no prime dividing all its integer values, takes positive prime values at infinitely many integer arguments. A missing sign qualification prevents treating this as a precise universal conjecture as written. The supplied open label and reconstruction claim require independent review.
\subsection{Sources}
\begin{itemize}
\item Local UnsolvedMath registry: supplied record with permanent ID \(30005267\), title ``Infinitely Many Prime Values of Integer Polynomials,'' exact source problem code \texttt{catalogue entry}, in the Oberwolfach Reports Open Problems collection. Project registry locations are \texttt{data/unsolvedmath/problems.json} and \texttt{catalog/problems.json}; this draft uses the supplied record without independently inspecting those files.
\item Supplied report citation: \emph{Analytic Number Theory (2023)}, Oberwolfach Reports, DOI \texttt{10.4171/owr/2022/50}, pp. 2895--2960. \url{https://doi.org/10.4171/owr/2022/50}. The record gives proposed year 2022 and no named proposer.
\item Supplied literature source: ``Oberwolfach Report 50.'' \url{https://ems.press/content/serial-article-files/46986}.
\item Supplied literature assessment, checked \textbf{2026-08-22}: the record labels the problem open and states that infinitude of prime values for irreducible polynomials of degree at least two with no fixed prime divisor is not known in general. This is supplied dated curation derived from the OWR material, not an independently verified literature assessment. The record itself requests independent verification of time-sensitive status and cited claims.
\item Supplied statement-verification metadata labels the statement \texttt{corrected\_verified} and says it was reconstructed and checked against the official Oberwolfach report. This is a claim in the supplied record, not verification performed for this draft.
\item Attribution in the original extracted passage: C\'ecile Dartyge, ``On the largest prime factor of quartic polynomial values: the cyclic and dihedral cases,'' joint work with James Maynard; an adjacent general result is attributed to Tenenbaum [6]. The supplied record does not identify the bibliographic entry [6] or establish that these authors proposed the prime-values question.
\end{itemize}
\end{document}
